% Отчет по лабораторной работе №6
% Дисциплина: Системы управления базами данных
% Оформление выполнено по шаблону ReportLab5.pdf
\documentclass[a4paper,12pt]{article}

\usepackage[T2A]{fontenc}
\usepackage[utf8]{inputenc}
\usepackage[russian]{babel}
\usepackage[a4paper,margin=2cm]{geometry}
\usepackage{graphicx}
\usepackage{adjustbox}

% Сансерифный набор, как в шаблоне отчета
\renewcommand{\familydefault}{\sfdefault}

% Нумерация подразделов вида 2.1., 2.2., ... (верхний уровень не печатается)
\setcounter{secnumdepth}{2}
\setcounter{section}{2}
\renewcommand{\thesubsection}{\thesection.\arabic{subsection}.}

% Сквозной счетчик рисунков (код и результаты нумеруются подряд)
\newcounter{shot}
\graphicspath{{./Screenshots/}}

% Вставка скриншота с подписью по канону шаблона:
%   Рисунок N -- Код запроса №M / Рисунок N -- Результат работы кода запроса №M
\newcommand{\codefig}[2]{%
  \stepcounter{shot}%
  \begin{center}
    \includegraphics[max width=0.95\textwidth,max height=0.72\textheight,keepaspectratio]{#1}%
    \par\vspace{6pt}
    Рисунок~\arabic{shot}~-- #2%
  \end{center}%
}

% Жирный заголовок-метка с новой строки (группы заданий, пункты а/б/в)
\newcommand{\grouplabel}[1]{\par\medskip\noindent\textbf{#1}\par\smallskip}

\begin{document}

% ============================ ТИТУЛЬНЫЙ ЛИСТ ============================
\begin{titlepage}
\centering
\vspace*{0.28\textheight}
{\large Системы управления базами данных\par}
\vspace{8pt}
{\large Лабораторная работа №6\par}
\vfill
\begin{flushright}
Выполнил: Сахань В.С.\\
ФИТ-242
\end{flushright}
\vspace*{0.05\textheight}
\end{titlepage}

% ============================ ХОД РАБОТЫ ============================
\subsection{Создание простых и материализованных представлений БД «Кинотеатр»}

Для проектируемой базы данных кинотеатра (CinemaDB) были сформированы три простых представления и два материализованных представления, обеспечивающих удобный доступ к операционной аналитике.

\grouplabel{1. Простые представления (Views)}

\textbf{Представление 1 (\texttt{vw\_active\_preorders}):} Отображает информацию об активных и выкупленных бронированиях с контактными данными клиентов.

\codefig{Code1.png}{Код запроса №1}
\codefig{CodeResult1.png}{Результат работы кода запроса №1}

\textbf{Представление 2 (\texttt{vw\_session\_details}):} Содержит расширенную информацию о сеансах, включая наименование фильма, жанр, название зала, базовую стоимость и количество проданных билетов.

\codefig{Code2.png}{Код запроса №2}
\codefig{CodeResult2.png}{Результат работы кода запроса №2}

\textbf{Представление 3 (\texttt{vw\_ticket\_sales}):} Детализирует каждый проданный билет с указанием сотрудника, оформившего продажу, и статуса брони.

\codefig{Code3.png}{Код запроса №3}
\codefig{CodeResult3.png}{Результат работы кода запроса №3}

\grouplabel{2. Материализованные представления (Materialized Views)}

\textbf{Материализованное представление 1 (\texttt{mv\_movie\_revenue\_stats}):} Фиксирует физически на диске сводную статистику по проданным билетам и общей выручке для каждого фильма.

\codefig{Code4.png}{Код запроса №4}
\codefig{CodeResult4.png}{Результат работы кода запроса №4}

\textbf{Материализованное представление 2 (\texttt{mv\_employee\_performance}):} Кэширует показатели эффективности работы кассиров (количество обслуженных билетов и суммарную выручку).

\codefig{Code5.png}{Код запроса №5}
\codefig{CodeResult5.png}{Результат работы кода запроса №5}

\subsection{Запросы к представлениям}

\grouplabel{а) Запрос на выборку, содержащий вложенный запрос}

Формулировка: Из материализованного представления \texttt{mv\_movie\_revenue\_stats} выбрать фильмы, суммарная выручка от продажи билетов на которые превышает среднюю выручку по всем фильмам кинотеатра.

\codefig{Code6.png}{Код запроса №6}
\codefig{CodeResult6.png}{Результат работы кода запроса №6}

\grouplabel{б) Запрос на выборку с группировкой данных и простой сортировкой}

Формулировка: По данным из представления \texttt{vw\_session\_details} сгруппировать сеансы по жанрам фильмов, вычислить общее количество запланированных сеансов и среднюю базовую стоимость билета в каждом жанре. Отсортировать результат по убыванию средней стоимости.

\codefig{Code7.png}{Код запроса №7}
\codefig{CodeResult7.png}{Результат работы кода запроса №7}

\grouplabel{в) Запрос на выборку с группировкой данных и многоуровневой сортировкой}

Формулировка: Сгруппировать данные о проданных билетах из представления \texttt{vw\_ticket\_sales} по должности сотрудника и статусу бронирования. Рассчитать количество билетов и общую сумму продаж. Выполнить многоуровневую сортировку: по роли сотрудника (в алфавитном порядке), а внутри одной роли --- по сумме продаж (по убыванию).

\codefig{Code8.png}{Код запроса №8}
\codefig{CodeResult8.png}{Результат работы кода запроса №8}

\subsection{Обновление материализованных представлений}

Так как материализованные представления хранят снимки данных, при изменении таблиц фактов (\texttt{Tickets}, \texttt{Sessions}) требуется выполнение команды обновления:

\codefig{Code9.png}{Код запроса №9}
\codefig{CodeResult9.png}{Результат работы кода запроса №9}

\end{document}
